\documentclass[10pt,twocolumn]{article}
\usepackage[margin=0.78in]{geometry}
\usepackage{amsmath,amssymb,booktabs,hyperref,microtype}
\usepackage[T1]{fontenc}

\title{Continuous Self-MTP Batching for Bandwidth-Bound\\
       Large-Language-Model Inference on Apple Silicon}
\author{Pierre Lamy}
\date{29 August 2026 --- working paper, not yet submitted}

\begin{document}
\maketitle

\begin{abstract}
Autoregressive decode on unified-memory systems is commonly limited by the
cost of streaming model weights for one token at a time.  Native multi-token
prediction (MTP) reduces the number of target-model rounds for a single
request, while continuous batching amortizes a target round across requests;
prior work has explored these dimensions both independently and jointly.  We
describe and characterize an MLX-specific realization of \emph{continuous
self-MTP batching}: a transactional decode engine that batches each request's
current token and its MTP-proposed continuation in one target forward.  For
$N$ homogeneous lanes and draft depth $k$, a cycle presents $M=(k+1)N$
verification positions while streaming the target weights once.  The design
adds per-lane residual verification, per-row ragged rollback, persistent
MTP-head state, dynamic lane admission, and prefix-cache composition.  In a
local MLX prototype, aggregate decode on
Qwen3.8-Flash-Next rose from 63.4 to 120.3 token/s between $N=1$ and $N=4$;
Qwen3.8-27B rose from 49.0 to 82.0 token/s.  These are session-trace
measurements from the source prototype, not measurements of the Rapid-MLX
port staged with this paper.  Flash is raw-log-attested and 27B is
transcript-attested; neither record is artifact-complete.  In the source
prototype, fixed cohorts and mid-flight joins for both families pass their
available hierarchical gates.  Rapid tip \texttt{d7cf5bbb} now contains a
default-off live coordinator and Qwen3.5 boundary joins/leaves; Rapid
Qwen4/Flash remains fixed-cohort and capability-refused for dynamic join.
\end{abstract}

\section{Motivation}
Plain decode authorizes one token after each target-model forward.  On Apple
Silicon, the parameters and KV state share one memory system, and large-model
decode is frequently bandwidth-bound.  A longer verification sequence can
therefore use otherwise under-filled compute while amortizing the dominant
weight stream.  Speculative decoding formalizes this exchange: a cheaper
proposal mechanism drafts tokens and the target verifies them without changing
the target distribution~\cite{leviathan2023}.  Trained prediction heads avoid
loading a separate draft model, as explored by Medusa and MTP-trained
models~\cite{cai2024,deepseek2024}.  Independently, iteration-level scheduling
changes batch membership between target rounds~\cite{yu2022}.

The systems problem is their composition.  A single-request self-MTP loop owns
target KV, recurrent state, an MTP-head cache, hidden-state handoff, random
draws, and pending tokens.  A continuous batch additionally permits requests
to arrive, finish, abort, hit a prefix cache, or become ineligible between
cycles.  Treating the MTP loop as an opaque per-request generator loses the
main opportunity: the target still streams once per lane.

\section{Continuous self-MTP}
Let lane $i$ propose $k_i$ tokens $d_{i,1:k_i}$.  Its target verification row
is
\begin{equation}
  v_i=[c_i,d_{i,1},\ldots,d_{i,k_i}],
\end{equation}
where $c_i$ is the committed token not yet represented in the target cache.
Rows are right padded to $W=1+\max_i k_i$ and evaluated together.  With fixed
depth $k$, the target sees
\begin{equation}
  M=(k+1)N
\end{equation}
logical verification positions per weight stream.  This is not a claim that
$M$ tokens are committed: lane $i$ commits $a_i$ accepted drafts plus one
target correction or bonus token.  Cycle yield is
\begin{equation}
  Y=\sum_{i=1}^{N}(a_i+1), \qquad
  \mathrm{TPS}=Y/t_{\mathrm{cycle}}.
\end{equation}

Draft generation is also batched.  Each native MTP-head step consumes the
per-lane hidden state and most recent token.  Lanes with a shorter remaining
budget contribute an invalid padded row rather than consuming RNG or cache
position.  Consequently, batch composition cannot change a lane's random
sequence.

For Flash-Next, the two recursive draft steps may share the QSA block
selection computed from each row's last valid first-step query.  The selection
is cycle-local: prepare/finalize padding inside the same proposal preserves it,
while commit, abort, state restore, trim, filter, extend, or membership change
clears it.  This is the Rapid adaptation of source commit
\texttt{393eddf}; it avoids both redundant indexer work and accidental reuse
across a geometry-changing boundary.  It does not by itself qualify dynamic
Flash joins.

\subsection{Verification and sampling}
Greedy verification accepts the longest prefix on which draft and target
argmax agree.  Sampled verification applies the same temperature, top-$p$,
top-$k$, and min-$p$ transform to both distributions, then uses the standard
acceptance probability
\begin{equation}
  \alpha(x)=\min\{1,p(x)/q(x)\}.
\end{equation}
On rejection it samples from the normalized residual
$[p-q]_+$.  Acceptance uniforms for lane $i$ have native shape $(k_i,)$ and
come from a lane-owned key.  They are never drawn as a padded $(N,W)$ tensor.
This preserves the lane's draw order across joins and leaves.  XTC remains
fail-closed until its stochastic transform can be shared exactly between
draft and target.  Block verification~\cite{sun2024} is compatible with the
transaction but is not required for the headline result.

Greedy digest qualification uses a separate hierarchy.  Level 1 compares the
batched engine at $B=1$ with incumbent single-lane self-MTP and requires exact
tokens and digest; no batch-shape tolerance is permitted.  Level 2 compares
each $B>1$ lane with its own batched-$B1$ execution and permits only the
declared BF16 batch-shape band.  The older single-lane-versus-$B>1$ comparison
is retained as an informative compatibility diagnostic.  Cache equality,
rollback, commit, and detach remain bit-exact gates, so distributional token
tolerance cannot mask transaction corruption.  This hierarchy follows source
commit \texttt{0995cbc}.  Rapid PR 12A encodes the offline classification
policy only: an external runner must calculate the numeric band, bind each
B$>$1 lane to its matching batched-B1 run, and persist the raw evidence.

\subsection{Propose--commit transaction}
Every cycle is a two-phase transaction:
\begin{enumerate}
  \item \textbf{Propose.} Snapshot membership epoch and lane order, advance
        the MTP head, verify the padded rows, and trim each target-cache row by
        $k_i-a_i$.
  \item \textbf{Commit.} After stop/EOS/length handling determines how many
        authorized tokens were actually delivered, trim any undelivered
        suffix, advance the canonical current token and hidden state, and
        publish the new lane state.
\end{enumerate}
Membership changes are forbidden while a proposal is open.  A nonterminal
lane must consume its complete authorized result.  A terminal lane may consume
only a prefix, but the target cache, recurrent slots, MTP cache, processor
history, and current-token boundary must all land on that same prefix.

This contract is essential for Qwen hybrid models.  Their logical cache is not
only transformer KV: it includes Gated DeltaNet convolution and recurrence,
PLE n-gram history, QSA raw/compressed index keys, and the MTP-head cache.  A
scalar rollback applied to only the transformer KV silently corrupts later
decode.

\subsection{Join, leave, and admission}
At a closed transaction boundary, leaving lanes are extracted before the old
batch is filtered.  Joining lanes arrive as canonical one-row packages and
are merged only after their capability and memory gates pass.  Each mutation
increments a membership epoch.  Model-family capability controls are
independent.  In the source prototype, 27B joins are token-exact and Flash-Next
joins pass the real-BF16 hierarchical gate distributionally: batched-B1 remains
exact, B$>$1 lanes---including the joined lane---stay inside the declared
batch-shape band, and cache/rollback contracts remain exact.  This source result
does not attest Rapid's distinct cache/runtime path.  Rapid wires Qwen3.5
dynamic membership behind default-off controls but keeps Qwen4/Flash
capability-refused pending Rapid-native BF16 qualification.

Admission is memory-aware and degrades monotonically:
\begin{equation}
  N \rightarrow k=2 \rightarrow k=1 \rightarrow
  \text{plain decode} \rightarrow \text{queue}.
\end{equation}
The Rapid-MLX planning layer reuses the server's existing allocator cap and
KV-footprint accounting instead of introducing a second memory model.
Decisions occur only between transactions, before widening a cache batch.
The current live path is opt-in, but its incremental draft-token cost is still
passed as zero, so depth-dependent memory degradation is not yet calibrated.

\section{Prefix-cache composition}
An exact prefix-cache record for self-MTP must represent one token boundary
across three objects: target cache $C_T$, persistent MTP cache $C_D$, and the
hidden state $h$ that predicts the current token.  A hit that restores only
$C_T$ either discards useful MTP work or resumes the draft head from an
inconsistent position.  The staged Rapid work therefore extends its existing
Flash-Next prefix-hit preservation; it does not replace Rapid's adaptive
prefix cache or radix index.

Prefix state is admitted only when offsets satisfy
\begin{equation}
  \mathrm{offset}(C_D)=\mathrm{offset}(C_T)-1,
\end{equation}
with a nonempty current token and a compatible model/configuration identity.
Trivial hits fail open to ordinary MTP initialization rather than disabling
MTP for the whole request.

\section{Implementation lineage}
The source implementation is preserved in four local commits:
\texttt{92576ced} (batched engine and Flash-Next gates),
\texttt{93d7aa9a} (27B/Qwen3.5 qualification and join tests),
\texttt{393eddf} (cycle-local QSA selection preservation), and
\texttt{0995cbc} (batched-B1 digest-oracle reframing), and
\texttt{1a0a2474} (compute-saturation admission ceiling and configurable
per-lane transient estimate).  Together they contain
an $N$-lane propose--commit core, a batch wrapper, server routing, per-lane RNG,
ragged cache lifecycle, admission, QSA draft reuse, and the final qualification
hierarchy.
The Rapid-MLX port targets \texttt{74652283} (release 0.13.1).  Its staged
series separates handoff correctness, pure transaction contracts, ragged
rollback, APC prepared state, the fixed-cohort core, family capabilities,
generation ownership, the MLX data plane, QSA selection sharing, transformed
residual sampling, planner routing, bounded qualification telemetry, and digest
oracle policy into reviewable commits.  Rapid commits \texttt{be0b3e8c} and
\texttt{d7cf5bbb} subsequently add wrapper membership and a combined runtime,
rollback, driver, scheduler-delivery, and Qwen3.5-turnover integration.  The
publication plan splits that broad final commit into narrower review layers.

Rapid already contains the Qwen4 architecture, vectorized QSA mask
construction, native one-layer MTP, batched compressed-key prefill, atomic
Qwen recurrent rollback, and MTP prefix-hit preservation.  Those features are
treated as prerequisites, not copied into the new stack.  The port focuses on
the mid-stream handoff invariant, batched prepared-state lifecycle, continuous
self-MTP coordinator, the prefix-state delta, and qualification/telemetry.

QSA kernel work remains outside the continuous self-MTP dependency chain.  The
Rapid-target side branch, historically named \texttt{qwen4-nax-side-stack},
actually rebases Rapid PR~\#2547's opt-in non-NAX Metal block-sparse prefill
kernel.  Its upstream M3 Ultra evidence reports a 10.7\% prefill gain at 32K
with unchanged decode; this local rebase did not independently rerun that
measurement.  Separately, the lab's MPP/NAX QSA prototype measured median
prefill gains of 16.0\%, 31.6\%, and 37.9\% at 16K, 32K, and 64K with flat
decode, preserved in committed per-arm JSON artifacts.  A session note records
that a later default-on production restart returned one request and the service
then disappeared without a new traceback.  No process-level artifact captured
that attempt, so it does not establish a kernel defect.  NAX remains
default-off pending an isolated, memory-observed rerun.  Neither kernel is part
of the headline self-MTP series.  NVMe PLE offload, shared-in-gather MoE, GDN normalization
fusion, and an untracked fused-expert kernel are also excluded because they
failed a gate or lack real-model validation.

\section{Evaluation}
Table~\ref{tab:throughput} reports aggregate decode rates recovered from the
source-development session trace.  $N=1$ is the optimized persistent self-MTP
loop, not plain autoregressive decode.  Four-lane values are aggregate service
throughput; approximate per-lane rates are 30.1 and 20.5 token/s.

\begin{table}[h]
\centering
\caption{Source-prototype aggregate decode throughput (token/s).}
\label{tab:throughput}
\begin{tabular}{lrrrr}
\toprule
Model & $N=1$ & $N=2$ & $N=4$ & $N=4/N=1$ \\
\midrule
Flash-Next & 63.4 & 88.4 & 120.3 & 1.90$\times$ \\
Qwen3.8-27B & 49.0 & 64.7 & 82.0 & 1.67$\times$ \\
\bottomrule
\end{tabular}
\end{table}

For 27B, a separate trace gate measured plain decode at 26.9 token/s and
single-lane self-MTP at 46.9 token/s (1.75$\times$), with 45 of 80 drafts
accepted and token-exact output.  Flash fixed-membership hardware validation
reported 9/9 bit-exact forced-cache comparisons and a 49-test suite.  The
available Rapid-MLX single-lane MTP benchmark at fixed $k=1$ reported
28.82--34.85 token/s across 92--32,732 prompt tokens and 76.41\% aggregate
acceptance; it was measured on an earlier Rapid commit and is background, not
a result for this port.

The table comes from Claude session
\texttt{1c3a4880-ce64-415d-9bb0-b06f7a29f73c}.  Flash-Next is preserved both
in transcript stdout and in the local raw log
\texttt{results/timing-batched-mtp.log} (SHA-256 prefix
\texttt{dc54cab994ec}); 27B is preserved in transcript stdout only because its
command piped through \texttt{grep}/\texttt{tail} without redirection.  A
tracked recovery manifest records the commands, trace lines, UUIDs, hashes,
and recovered rows.  Neither record is an artifact-complete benchmark package:
the final battery lacks a repository-tracked original JSON/JSONL/CSV containing
immutable revisions, per-sample timings, thermal/interleave order, digests,
memory, and acceptance.  We therefore label Flash raw-log-attested and 27B
transcript-attested, pending a machine-readable rerun.

The initial Rapid PRs 1--13 staging effort was model-free.  A later limited
smoke exercised the real Qwen3.8-27B target trunk with an explicitly test-only
randomly initialized MTP head.  Acceptance was zero at every cohort size; the
observed ladder rose from 18.4 token/s at $N=1$ to 87.3 at $N=16$
(4.74$\times$).  That run exercised real-Metal runtime assembly, correction-token
delivery, a mid-flight join, per-lane leave, and cleanup, but it has no
checked-in raw artifact.  A checked-in recovery CSV binds every row to the
session trace, random-head warning, and zero-acceptance boundary; it is not a
substitute for raw output.  The run does not qualify trained-head acceptance,
accepted-draft commit behavior, production throughput, or self-MTP speedup.
The Rapid stack is therefore not claimed merge-ready or
performance-qualified.  Its acceptance matrix must use
immutable checkpoints, process isolation, interleaved thermal controls,
prompt/seed manifests, per-lane and aggregate rates, acceptance, active and
peak memory, cache equality, abort/EOS behavior, APC hits, and join/leave
discriminators.

\section{Excluded work}
The staged series deliberately excludes the following paths:
\begin{itemize}
  \item Rapid Qwen4/Flash dynamic joins, which remain capability-refused until
        the distinct Rapid path passes the real-BF16 hierarchical gate;
  \item NVMe PLE offload, whose swap abort gate failed;
  \item \texttt{moe\_shared\_in\_gather}, which diverged in logprobs and
        output digest;
  \item GDN q/k normalization fusion, which regressed decode;
  \item the untracked fused-expert kernel, a default-off candidate without
        real-model validation;
  \item the old thread-crossing RNG wiring---per-lane keys remain owned by the
        generation thread; and
  \item quantized or windowed batched-MTP caches until their ragged rollback
        contracts have independent tests.
\end{itemize}

\section{Limitations and open work}
First, the live scheduler removes requests from its queue before driver
creation or join is known to succeed; publication requires construct-first or
exception restoration plus fault injection.  Second, APC prepared state exists
as a planner contract but is not consumed by the live runtime.  Third, the live
route is greedy-only even though transformed verification exists offline.
Fourth, live admission currently passes zero incremental draft-token memory
cost, so depth-dependent degradation needs calibration.  Fifth, aggregate
throughput trades against per-lane inter-token latency and must be evaluated
under arrival processes.  Heterogeneous depth and quantized/windowed caches
need independent proofs.  Finally, source Flash join evidence does not qualify
Rapid Flash, and results at one unified-memory operating point do not establish
portability across model sizes, quantization, context lengths, or Apple SoCs.

\section{Conclusion}
Continuous self-MTP batching treats speculative decode as an iteration-level
transaction rather than a private per-request generator.  It composes the two
available batching axes---tokens within a speculative block and requests
within a service cohort---while preserving per-lane distribution, rollback,
and cache ownership.  The source results show substantial aggregate-throughput
headroom.  The staged Rapid integration separates the historical planner,
wrapper membership, runtime/rollback, fixed-cohort live delivery, and Qwen3.5
turnover into review layers.  Qwen3.5 dynamic membership is wired but not
trained-head-qualified; Qwen4/Flash remains refused.  This lets later hardware
qualification promote each capability without weakening the fail-closed
boundary.

\section{Prior art and novelty boundary}
The broad concepts of self-speculative MTP, batched speculative verification,
and continuously scheduled speculative serving are established prior art.
Meta's MTP work used a model's own prediction heads for self-speculative
decoding and explicitly supported heterogeneous batch sizes~\cite{metamtp2024}.
BASS jointly parallelized the request and draft-token dimensions, including
per-sequence acceptance lengths, ragged KV handling, and sharing model-weight
reads across the batch~\cite{bass2024}.  TurboSpec incorporated
speculative proposal and collective target verification into an online
continuous-batching system with adaptive proposal lengths~\cite{smartspec2024}.
DeepSeek-V3 established recurrent native MTP as a speculative drafter, although
its report did not specify the continuous-serving cache lifecycle
\cite{deepseek2024}.

Production-oriented open-source systems further narrow the claim.  vLLM added
native DeepSeek MTP to its continuously scheduled online server and tested
multi-request execution and preemption~\cite{vllmmtp}.  SGLang documents native
MTP at substantial concurrency and experimental draft/verification overlap
\cite{sglangmtp}.  TensorRT-LLM documents recurrent MTP-head reuse, a distinct
MTP KV state, rejected-token eviction and target-cache rewind, and integration
with its overlap scheduler~\cite{trtmtp}.  AngelSpec treats shared batch-level
verification as a resource allocated across concurrent requests; its closest
cross-request mechanism is evaluated separately from its recurrent-MTP arm,
so it is strong adjacent rather than evidence of the exact construction here
\cite{angelspec2026}.

The closest MLX precedent found in this sweep is AirRunner's batched native-MTP
branch.  Commit \texttt{79a8367} implements $B>1$ MTP with per-sequence
acceptance, \texttt{GenerationBatch.extend()} and filtering, and selective SSM
and KV rollback; \texttt{99d0cc8} follows by committing accepted-token MTP
cache state in a batched forward~\cite{airrunnermtp}.  The commit objects carry
May 2026 author dates and August 2026 rebased committer dates, both preceding
the present 29 August prototype.  This is direct prior art for generic claims
about batched native MTP, batch resize/filter primitives, or selective
hybrid-cache rollback.  Its published cycle drafts one MTP token and verifies $M=2N$;
the present engine targets recursive depth two and $M=3N$, and extends the
transaction to Qwen's GDN, PLE, QSA, APC, lane-owned RNG, and admission state.

The \texttt{mlxcel} runtime independently added $B>1$ MTP/DFlash server cohorts
in May 2026 and later admitted ragged prompt lengths~\cite{mlxcelmtp}.  Its
current scheduler documentation states that each $B>1$ speculative burst is a
compatible cohort fixed at dispatch and run to completion inside one scheduler
tick.  It is therefore direct Apple/MLX prior art for batched MTP, but not for
admitting replacements at transaction boundaries inside an already-live
$B>1$ speculative service cohort.

Accordingly, this paper does \emph{not} claim the invention or first use of
batched MTP, continuous speculative serving, per-lane rollback, or their broad
combination.  Its contribution is the implementation and characterization of
a narrower conjunction: depth-two $M=3N$ execution for Qwen hybrid state in
MLX; a single propose--commit boundary spanning KV, recurrence, PLE, QSA, and
MTP-head state; lane-local transformed-sampling RNG and APC ownership; and, in
the source prototype, dual memory/compute admission validated at distinct dense
and MoE operating points.  This bounded primary-source sweep found no
disclosure of that full
conjunction.  ``No exact prior found'' is a scoped research result, not proof
of absence.

\begin{thebibliography}{99}
\bibitem{leviathan2023}
Y. Leviathan, M. Kalman, and Y. Matias.
\newblock Fast inference from transformers via speculative decoding.
\newblock \emph{ICML}, 2023. \url{https://arxiv.org/abs/2211.17192}.

\bibitem{cai2024}
T. Cai et al.
\newblock Medusa: Simple LLM inference acceleration framework with multiple
decoding heads.
\newblock \emph{ICML}, 2024. \url{https://arxiv.org/abs/2401.10774}.

\bibitem{deepseek2024}
DeepSeek-AI et al.
\newblock DeepSeek-V3 technical report.
\newblock 2024. \url{https://arxiv.org/abs/2412.19437}.

\bibitem{sun2024}
Z. Sun et al.
\newblock Block verification accelerates speculative decoding.
\newblock 2024. \url{https://arxiv.org/abs/2403.10444}.

\bibitem{yu2022}
G.-I. Yu et al.
\newblock Orca: A distributed serving system for transformer-based generative
models.
\newblock \emph{OSDI}, 2022.
\url{https://www.usenix.org/conference/osdi22/presentation/yu}.

\bibitem{metamtp2024}
A. Gloeckle et al.
\newblock Better \& faster large language models via multi-token prediction.
\newblock 2024. \url{https://arxiv.org/abs/2404.19737}.

\bibitem{bass2024}
H. Qian et al.
\newblock BASS: Batched attention-optimized speculative sampling.
\newblock 2024. \url{https://arxiv.org/abs/2404.15778}.

\bibitem{smartspec2024}
X. Liu et al.
\newblock TurboSpec: Closed-loop speculation control system for optimizing
LLM serving goodput.
\newblock 2024. \url{https://arxiv.org/abs/2406.14066}.

\bibitem{vllmmtp}
vLLM contributors.
\newblock DeepSeek native-MTP speculative decoding, pull request 12755.
\newblock 2025.
\url{https://github.com/vllm-project/vllm/pull/12755}.

\bibitem{sglangmtp}
SGLang contributors.
\newblock DeepSeek-V3 and native-MTP serving documentation.
\newblock \url{https://github.com/sgl-project/sglang/blob/main/docs/basic_usage/deepseek_v3.md}.

\bibitem{trtmtp}
NVIDIA TensorRT-LLM contributors.
\newblock DeepSeek-R1 MTP implementation and optimization.
\newblock \url{https://github.com/NVIDIA/TensorRT-LLM/blob/main/docs/source/blogs/tech_blog/blog02_DeepSeek_R1_MTP_Implementation_and_Optimization.md}.

\bibitem{angelspec2026}
H. Liu et al.
\newblock AngelSpec: Towards real-world high performance inference with
speculative decoding.
\newblock 2026. \url{https://arxiv.org/abs/2607.25852}.

\bibitem{airrunnermtp}
AirRunner.
\newblock Batched speculative decoding ($B>1$) with selective rollback;
accepted-token MTP-cache follow-up.
\newblock 2026.
\url{https://github.com/AirRunner/mlx-lm/commit/79a836750bd5b363cc1ec9a24df85741376afe5e};
\url{https://github.com/AirRunner/mlx-lm/commit/99d0cc8b090e5e8322534fe31f695b607bd1754e}.

\bibitem{mlxcelmtp}
lablup \texttt{mlxcel} contributors.
\newblock Batched MTP/DFlash server dispatch and continuous-batching design.
\newblock 2026.
\url{https://github.com/lablup/mlxcel/commit/6a816947e11882b5f174156c7c483acfd00ca409};
\url{https://github.com/lablup/mlxcel/blob/main/docs/CONTINUOUS_BATCHING.md}.
\end{thebibliography}

\end{document}
